%% Abstract text
%% All the details (name, title, etc.) on the abstract page appear as specified
%% above. Add your abstarct text with paragraphs here to have paragraphs in the
%% visible abstract page. Nonetheless, write the abstarct text without
%% paragraphs in the macro \thesismacro so that it is added to the metadata as
%% well.
%%
\begin{abstractpage}[english]

\todo{Write the abstract after everything settled.}

  % The abstract is a short description of the essential contents of the thesis,
  % usually in one paragraph: what was studied and how and what were the main
  % findings.

  % For a Finnish thesis, the abstract should be written in both Finnish and
  % English; for a Swedish thesis, in Swedish and English. The abstracts for
  % English theses written by Finnish or Swedish speakers should be written in
  % English and either in Finnish or in Swedish, depending on the student’s
  % language of basic education. Students educated in languages other than Finnish
  % or Swedish write the abstract only in English. Students may include a second
  % or third abstract in their native language, if they wish.

  % The abstract text of this thesis is written on the readable abstract page as
  % well as into the pdf file's metadata via the \verb+\thesisabstract+ macro
  % (see comment in this \TeX{} file above). Write here the text that goes onto
  % the readable abstract page. You can have special characters, linebreaks, and
  % paragraphs here. Otherwise, this abstract text must be identical to the
  % metadata abstract text.
  
  % If your abstract does not contain special characters and it does not require
  % paragraphs, you may take advantage of the \verb+\abstracttext+ macro (see the
  % comment in this \TeX{} file below).
\end{abstractpage}

%% The text in the \thesisabstract macro is stored in the macro \abstractext, so
%% you can use the text metadata abstract directly as follows:
%%
%\begin{abstractpage}[english]
%	\abstracttext{}
%\end{abstractpage}

%% Force a new page so that the possible Finnish or Swedish abstract does not
%% begin on the same page
%%
% \newpage
%%
%% Abstract in Finnish. Delete if you don't need it. 
%%
%% Respecify those fields that differ from the earlier specification or simply
%% respecify all fields.
% \thesistitle{Opinnäyteen otsikko}
% \thesissubtitle{Opinnäytteen mahdollinen alaotsikko}
% \supervisor{Prof.\ Pirjo Professori}
% \advisor{Yliopisto-opettaja Juho Vepsäläinen}
% \degreeprogram{Elektroniikka ja sähkötekniikka}
% \major{Sopiva pääaine}
% \collaborativepartner{Yhtiön tai laitoksen nimi (tarvittaessa)}
% \date{30.6.2025}
% %% The keywords need not be separated by \spc now.
% \keywords{Vastus, resistanssi, lämpötila}
% %% Abstract text
% \begin{abstractpage}[finnish]
% Tiivistelmä on lyhyt kuvaus työn keskeisestä sisällöstä usein yhtenä kappaleena:
% mitä tutkittiin ja miten sekä mitkä olivat tärkeimmät tulokset. Suomenkielisen
% opinnäytteen tiivistelmä kirjoitetaan suomeksi ja englanniksi ja ruotsinkielisen
% vastaavasti ruotsiksi ja englanniksi. Suomen- tai ruotsinkielisten
% opiskelijoiden, joiden opinnäytteen kieli on englanti, tulee kirjoittaa
% tiivistelmänsä englanniksi ja koulusivistyskielellään. Muiden kuin
% koulusivistyskieleltään suomen- tai ruotsinkielisten tulee kirjoittaa
% tiivistelmänsä vain englanniksi. Opiskelija voi halutessaan lisätä
% opinnäytteeseensä toisen tai kolmannen tiivistelmän omalla äidinkielellään.
% Tämän opinnäytteen tiivistelmäteksti kirjoitetaan opinnäytteen luettavan osan
% lomakkeen lisäksi myös pdf-tiedoston metadataan. Kirjoita tähän metadataan
% kirjoitettavaa teksti. Metadatatekstissa ei saa olla erikoismerkkejä,
% rivinvaiho- tai kappaleenjakomerkkiä, joten näitä merkkeja ei saa käyttää tässä.
% Jos tiivistelmäsi ei sisällä erikoimerkkejä eikä kaipaa kappaleenjakoa, voit
% hyödynttää makroa abstracttext luodessasi lomakkeen tiivistelmää (katso
% kommentti tässä TeX-tiedostossa alla). Metadatatiivistelmatekstin on muuten 
% oltava sama kuin lomakkeessa oleva teksti.

% \end{abstractpage}

% %% Force new page so that the Swedish abstract starts from a new page
% \newpage

% %% Swedish abstract. Delete it if you don't need it. 
% %% 
% %% Respecify those fields that differ from the earlier specification or simply
% %% respecify all fields.
% \thesistitle{Arbetets titel}
% \supervisor{Prof.\ Pirjo Professori}
% \advisor{TkD Alan Advisor} %
% \advisor{DI Elsa Expert}
% \degreeprogram{Electronik och electroteknik}
% \collaborativepartner{Company or institute name in Swedish (if relevant)}
% %\date{30.6.2025}
% %% Abstract keywords
% \keywords{Nyckelord p\aa{} svenska, temperatur}
% %% Abstract text
% \begin{abstractpage}[swedish]
% Sammandraget är en kort beskrivning av arbetets centrala innehåll: vad 
% undersöktes, hur undersöktes det och vilka var de viktigaste resultaten?

% I lärdomsprov som skrivs på svenska skrivs sammandraget på svenska och engelska, 
% på motsvarande sätt skrivs sammandraget på finska och engelska i lärdomsprov på 
% finska. Finsk- eller svenskspråkiga studerande som skriver sitt lärdomsprov på 
% engelska ska skriva sammandraget på engelska och på sitt skolutbildningsspråk. 
% Studerande vars skolutbildningsspråk inte är svenska eller finska skriver 
% sammandraget endast på engelska. Den studerande kan om hen så önskar lägga till 
% ett andra eller tredje sammandrag på sitt eget modersmål. Sammandraget fungerar 
% då ofta som mognadsprov och bör i så fall vara minst 300 ord långt. Information 
% om mognadsprov på svenska finns på MyCourses:\\
% \url{https://mycourses.aalto.fi/course/view.php?id=26872}.
% \end{abstractpage}
